\section{Flavour decomposition}
\label{app:flavour}

% The timing table is defined here to control its placement.
\begin{table}[t]
\centering
\caption{Evaluation time for the two complete Type II coefficients at the benchmark parameters, $c=\TimingCos$ and \TimingPrecision-digit working precision, for the unexpanded remainders (exact) and the series. First calls start with empty caches; repeated times are medians of \TimingRepeats{} evaluations at the same point. Starting the kernel and loading the files are excluded.}
\label{tab:timing}
\begin{ruledtabular}
\footnotesize
\setlength{\tabcolsep}{2.6pt}
\begin{tabular}{lcccccc}
& \multicolumn{3}{c}{first call [s]} & \multicolumn{3}{c}{repeated [s]} \\
$\bA$ & exact & series & ratio & exact & series & ratio \\
\hline
$1/10$ & 56.04 & 11.61 & 4.83 & 0.3422 & 0.0178 & 19.18 \\
$1/5$ & 45.68 & 11.64 & 3.92 & 0.3403 & 0.0179 & 18.99 \\
$3/10$ & 44.64 & 11.64 & 3.84 & 0.3485 & 0.0181 & 19.22 \\
\end{tabular}
\end{ruledtabular}
\end{table}


To compare the sizes of the flavour contributions, I evaluate the decomposition of
Eq.~\eqref{eq:nf-decomposition} at $\nf=\BenchNf$:
\begin{align}
 \Mfull{i}(\BenchNf)
 &=\Mflav{i}{0}
   +\BenchNf\,\Mflav{i}{1}
   +\BenchNf^2\,\Mflav{i}{2}
 \nonumber\\
 &\equiv a_0+a_1+a_2,
 \qquad i=1,2.
 \label{eq:nf5-sum}
\end{align}
The index $i$ is suppressed in $a_r$. The factors $\BenchNf$ and $\BenchNf^2$
are included in $a_1$ and $a_2$, respectively, so these quantities are the three
terms entering the complete remainder, rather than the unweighted flavour coefficients.

Figure~\ref{fig:flavour-components} shows the magnitude of each term divided by
the magnitude of the complete remainder, $|a_r|/|\Mfull{i}(\BenchNf)|$, at
$c=\BenchCos$, $\mA=\muR=\BenchMA\,\mathrm{GeV}$ and $\Nc=\BenchNc$. Panels~(a) and~(b)
correspond to $\mathcal M_1$ and $\mathcal M_2$, respectively. Since the terms are
complex, their magnitudes do not add to the magnitude of their sum.
The plotted ratios therefore need not add to unity. A value above
the horizontal line at unity means that an individual term is larger
in magnitude than the complete remainder and is partly cancelled
by the other terms.

The term independent of $\nf$ is
the largest contribution to both coefficients over the plotted
range. The term linear in $\nf$ is also sizeable. For
$\mathcal M_1$, it partly cancels the $\nf^0$ contribution.
The overall cancellation among the three terms can be quantified
by
\[
 1-\frac{|a_0+a_1+a_2|}
          {|a_0|+|a_1|+|a_2|}.
\]
This quantity is zero when the terms add with the same phase
and increases as their sum is reduced by cancellation.
For $\mathcal M_1$, it lies between \PoneCancellationMin{}
and \PoneCancellationMax{} over the velocities shown.
These percentages refer to the reduction in magnitude relative
to $|a_0|+|a_1|+|a_2|$, not to a reduction in the cross section.

The relative signs of the flavour-dependent and gluonic
terms can also be seen in the leading power coefficients.
The combinations $\beta_0=(11\Nc-2\nf)/3$,
$\Nc a_{\mathrm A}-\nf a_{\mathrm F}$ and $4\Nc-\nf$ enter
$\widehat A_{12}$, $\widehat A_{11}$ and
$\widehat A_{21}$, respectively [Eqs.~\eqref{eq:leading-coefficients}
and \eqref{eq:M1-L1}]. They exhibit the relative minus
signs in the logarithmic terms; the size of the cancellation
in the complete remainder also depends on their phases and
on the powers of $\mathcal L$ that they multiply.

Finally, the $\nf^2$ contribution to $\mathcal M_2$ starts at
$\bA^2(1-c^2)$, as shown in Eq.~\eqref{eq:nf2-leading-two}.
This suppression explains why its relative contribution in
panel~(b) becomes small near threshold.
